\subsection{一致结构的比较}

第 1 目的命题 2 表明，我们可以取一致连续映射作为一致结构的态射（《集合论》，第四章，§ 2，第 1 目）；今后我们总假定态射就是这样选取的。按照一般的定义（《集合论》，第四章，§ 2，第 2 目），这使我们可以对给定集合 X 上的一切一致结构组成的集定义一个序关系：

定义 2 若 \(\mathfrak{U}_1\) 与 \(\mathfrak{U}_2\) 是同一集合 \(\mathbf{X}\) 上的两个一致结构，则称 \(\mathfrak{U}_1\) 比 \(\mathfrak{U}_2\) 细（而 \(\mathfrak{U}_2\) 比 \(\mathfrak{U}_1\) 粗），如果记 \(\mathbf{X}_i\) 为集合 \(\mathbf{X}\) 赋以一致结构 \(\mathfrak{U}_i\) 后所得的空间（ \(i = 1,2\)），恒等映射 \(\mathbf{X}_1 \to \mathbf{X}_2\) 一致连续。

若 \(\mathfrak{U}_1\) 比 \(\mathfrak{U}_2\) 细且异于 \(\mathfrak{U}_2\)，则称 \(\mathfrak{U}_1\) 严格地细于 \(\mathfrak{U}_2\)（而 \(\mathfrak{U}_2\) 严格地粗于 \(\mathfrak{U}_1\)）。

两个一致结构称为可比较的，如果其中一个比另一个细。

例。在集合 X 上一切一致结构组成的有序集中，离散一致结构是最细的，而最粗的一致结构是其一致邻域集由单个元素 \(X \times X\) 组成的那一个。

下述命题是第 1 目定义 1 的直接推论：

命题 3 若 \(\mathfrak{U}_1\) 与 \(\mathfrak{U}_2\) 是集合 \(X\) 上的两个一致结构，则 \(\mathfrak{U}_1\) 比 \(\mathfrak{U}_2\) 细，当且仅当 \(\mathfrak{U}_2\) 的每个一致邻域都是 \(\mathfrak{U}_1\) 的一致邻域。

推论 设 \(\mathfrak{U}_1\) 与 \(\mathfrak{U}_2\) 是集合 \(X\) 上的两个一致结构，且设 \(\mathfrak{U}_1\) 比 \(\mathfrak{U}_2\) 细；则 \(\mathfrak{U}_1\) 诱导的拓扑比 \(\mathfrak{U}_2\) 诱导的拓扑细。

这由拓扑按邻域的比较（第一章 § 2 第 2 目命题 3）立即得出。

注。1）可能发生这样的情形：一致结构 \(\mathcal{U}_1\) 严格地细于一致结构 \(\mathcal{U}_2\)，但二者诱导的拓扑相同。下面的例子表明这一点：

设 \(\mathbf{X}\) 是非空集合。对每个有限分拆 \(\varpi = (\mathrm{A}_i)_{1\leqslant i\leqslant n}\)（\(\mathbf{X}\) 的），令 \(\mathrm{V}_{\overline{\omega}}\) 表示

\[
\bigcup_ {i} \mathrm{A} _ {i} \times \mathrm{A} _ {i}.
\]

于是诸集合 \(V_{\overline{w}}\) 构成某个一致结构 \(\mathcal{U}\)（\(X\) 上的）的一致邻域基本系。因为若 \(\varpi\) 是 \(X\) 的任一有限分拆，我们有 \(\Delta \subset V_{\overline{w}}\) 且 \(V_{\overline{w}} \circ V_{\overline{g}} = \overline{V}_{\overline{w}}^{1} = V_{\overline{w}}\)（§ 1 第 1 目例 2）；又若 \(\varpi' = (B_j)\) 与 \(\varpi'' = (C_k)\) 是 \(X\) 的两个有限分拆，则诸集合 \(B_j \cap C_k\) 中非空者构成一个有限分拆 \(\varpi\)（\(X\) 的），并且 \(V_{\overline{w}} \subset V_{\overline{w}'} \cap V_{\overline{w}''}\)。称 \(\mathcal{U}\) 为 \(X\) 上有限分拆的一致结构。由 \(\mathcal{U}\) 诱导的拓扑是离散拓扑，因为对每个 \(x \in X\)，集合 \(\{x\}\) 与 \(C\{x\}\) 构成 \(X\) 的一个有限分拆。然而，若 \(X\) 无限，则显然 \(\mathcal{U}\) 严格地粗于离散一致结构。2）若 \(f: X \to X'\) 是一致连续映射，则 \(f\) 在把 \(X\) 的一致结构换成较细的一致结构、把 \(X'\) 的一致结构换成较粗的一致结构之后仍一致连续（第 1 目命题 2）。换句话说，\(X\) 的一致结构越细且 \(X'\) 的一致结构越粗，\(X\) 到 \(X'\) 的一致连续映射就越多。

